\documentclass[10pt,a4paper]{amsart}
\usepackage[margin=19mm]{geometry}
\usepackage{amsmath,amssymb,mathtools}
\usepackage[scr=rsfs,cal=euler]{mathalfa}
\usepackage{xfrac}
\usepackage[unicode,hidelinks,bookmarks=false]{hyperref}
\newcommand{\ie}{i.e.\ }
\newcommand{\cf}{cf.\ }
\newcommand{\resp}{respectively\ }
\newcommand{\Subsection}{\subsection}
\providecommand{\nobreakdash}{\nobreak\hbox{-}}
\setlength{\emergencystretch}{2em}
\multlinegap0pt
\usepackage{bm}
\usepackage{bbm}
\usepackage{dsfont}
\usepackage{xspace}
\usepackage{cancel}
\usepackage{mathtools}
\usepackage{amsthm}
\makeatletter
\@ifundefined{theorem}{\newtheorem{theorem}{Theorem}}{}
\@ifundefined{claim}{\newtheorem{claim}[theorem]{Claim}}{}
\makeatother
\renewcommand{\d}{{\downarrow}}
\renewcommand{\u}{{\uparrow}}
\title[A topos for \'etale-finite Heyting algebras]{A topos for \'etale-finite Heyting algebras}
\author{Marco Abbadini, Rodrigo Nicolau Almeida, Igor Arrieta}
\date{Source: arXiv:2606.03861v2 (CC BY 4.0)}
\begin{document}
\maketitle

\noindent\emph{Source and licence.} A topos for \'etale-finite Heyting algebras. Version arXiv:2606.03861v2 (CC BY 4.0), distributed under the Creative Commons Attribution 4.0 International licence, as recorded in the arXiv licence metadata for this exact version and shown on the abstract page https://arxiv.org/abs/2606.03861v2. Full rights notice: licenses/arxiv-2606.03861-CC-BY-4.0.txt.\par\smallskip

\noindent\textit{Editorial scope.} Counted unit: the Claim whose source label is \texttt{claim:finite-clopen-partition-separating-principal-upsets} in the article (source0.tex lines 2491--2498, source label \texttt{claim:finite-clopen-partition-separating-principal-upsets}), delivered together with the article's own complete original proof of it (source0.tex lines 2500--2549, one \texttt{proof} environment of 50 lines). No mathematics is written, repaired or re-derived: statement and proof are the article's own text line for line. Required context retained uncounted: none. Every definition, display and label of those lines is retained unchanged. Omitted from this package and not redistributed: 1--2490, 2499--2499, 2550--2817. Nothing inside a retained excerpt is omitted or altered: every retained body is delivered exactly as the article's lines are written, so no omission receipt of any kind accompanies this package. Citations: the retained excerpts carry no citation command at all, so no bibliography entry of the article is transcribed and no reference is redistributed. Reference adaptation: none was needed and none was made; every reference inside the delivered unit resolves inside it, so no unresolved reference or citation is produced. Printed slips kept literal: no printed word, symbol, hypothesis or number of the counted statement or of its proof was altered. Layout and class: the article's own class is \texttt{amsart} with packages including babel, microtype, fullpage, mathtools, amsthm, amsfonts, amssymb, stmaryrd, xcolor, graphicx, tikz-cd, todonotes, hyperref, cleveref; the wrapper is the lane's pinned \texttt{amsart} editorial wrapper at 10pt on A4, which restates the theorem-like environments the retained text needs and adds the article's own macro definitions that the retained text needs (\texttt{\textbackslash d}, \texttt{\textbackslash u}), with extra wrapper packages (bm, bbm, dsfont, xspace, cancel, mathtools). The delivered numbering is the unit's own. Assets: no retained span contains a figure environment, a graphics-inclusion command, a PDF-page inclusion, a table or a TikZ picture; no image file is copied into this package, the package declares no asset dependency and no image file is redistributed with it. Licence: version arXiv:2606.03861v2 of this article is distributed under the Creative Commons Attribution 4.0 International licence, as recorded in the arXiv licence metadata for this exact version and shown on the abstract page https://arxiv.org/abs/2606.03861v2. A full rights notice is delivered at licenses/arxiv-2606.03861-CC-BY-4.0.txt. Characters: every retained excerpt is pure ASCII, so the delivered engine, XeLaTeX with its default Latin Modern text font, reports no missing character.
    \begin{claim} \label{claim:finite-clopen-partition-separating-principal-upsets}
        There is a finite clopen partition
        \[
            B_1,\dots,B_r
        \]
        of $X$ such that, for every $x\in X$, distinct points of ${\u}x$ lie in distinct
        blocks of the partition.
    \end{claim}

    \begin{proof}[Proof of the claim]
        Let $x\in X$, and write $n\coloneqq \lvert{\u}x\rvert$. Since the underlying topological space of a Priestley space is a Stone space, we may choose pairwise
        disjoint clopen neighbourhoods
        \[
            U^x_y\qquad(y\in{\u}x)
        \]
        such that $y\in U^x_y$ for each $y\in{\u}x$. Define
        \[
            \mathcal{O}_x
            \coloneqq
            X_n
            \cap
            \Box\left(\bigcup_{y\in{\u}x}U^x_y\right)
            \cap
            \bigcap_{y\in{\u}x}{\d}U^x_y.
        \]
        Then $\mathcal{O}_x$ is clopen. Indeed, $X_n$ is clopen, the set
        $\bigcup_{y\in{\u}x}U^x_y$ is clopen, and both $\Box W$ and ${\d}W$ are clopen
        whenever $W$ is clopen in an Esakia space.

        Moreover, $x\in\mathcal{O}_x$. Thus the family $(\mathcal{O}_x)_{x\in X}$
        covers $X$, and hence, by compactness, there are $x^1,\dots,x^m\in X$ such that
        \[
            X=\mathcal{O}_{x^1}\cup\cdots\cup\mathcal{O}_{x^m}.
        \]

        Collect the finitely many clopen subsets
        \[
            U^{x^a}_y
            \qquad
            (1\leq a\leq m,\ y\in{\u}x^a),
        \]
        and let $B_1,\dots,B_r$ be the finite clopen partition of $X$ generated by them.

        We show that this partition has the desired property. Let $z\in X$, and choose
        $a$ such that $z\in\mathcal{O}_{x^a}$. Set
        \[
            n\coloneqq \lvert{\u}x^a\rvert.
        \]
        Since $z\in X_n$, we have $\lvert{\u}z\rvert=n$. Moreover, by the definition of
        $\mathcal{O}_{x^a}$, the set ${\u}z$ is contained in
        \[
            \bigcup_{y\in{\u}x^a}U^{x^a}_y
        \]
        and meets each of the pairwise disjoint clopens $U^{x^a}_y$. Since both ${\u}z$
        and ${\u}x^a$ have cardinality $n$, the set ${\u}z$ meets each $U^{x^a}_y$ in
        exactly one point. Hence, any two distinct points of ${\u}z$ are separated by one of
        the chosen clopens, and therefore lie in distinct blocks of the generated partition.
        This proves the claim.
    \end{proof}

\end{document}
